\chapter*{Abstract}
\thispagestyle{empty}
Active Galactic Nuclei (AGN) are intensely luminous objects powered by supermassive black holes at the centres of galaxies. Their radiant emissions can often overpower the light from their host galaxy, making it challenging to accurately measure galaxy properties such as star formation and redshift. Colour-colour diagnostics like the UVJ and ugr diagrams are widely utilised tools to estimate these properties using galaxy photometry. However, the contribution of AGN can significantly alter a galaxy's position within these diagnostics, potentially leading to misclassification and inaccuracies in galaxy characterisation.\\\\By utilising both theoretical models and observational data from the ZFOURGE Survey, we investigate how AGN affect galaxy properties estimated using colour-colour diagnostics. To do this, we generated theoretical composite spectral energy distributions (SEDs) with varying AGN contributions by integrating theoretical AGN models with a suite of semi-empirical galaxy templates. We then used synthetic filters to generate photometry from these SEDs. Analysis of synthetic photometry from these SEDs revealed that unobscured AGN (Type 1 AGN), even at low contributions, induce a significant shift towards bluer colours in both UVJ and ugr diagrams, up to a mean vector offset of 0.52 and 1.52 dex, respectively. This shift can lead to the misclassification of quiescent galaxies as star-forming in the UVJ space. It may also lead to an underestimation of redshift when using the ugr diagram.
In contrast, AGN obscured by their environment (Type 2 AGN) exhibited minimal impact on these colour spaces. Observational data from the ZFOURGE survey supported these findings. When AGN light was removed from a source using SED fitting and decomposition, galaxies exhibited a shift in the UVJ colour space, effectively reversing the impact of AGN. It is suggested that discrepancies in the shift between the theoretical and observational data may be attributed to an intermediate-type AGN, which requires further investigation.\\\\ 
This research highlights the potential for AGN to contaminate colour-colour diagrams, leading to galaxy misclassification and incorrect photometric redshift determination in widely used colour diagnostics. This underscores the importance of considering AGN contributions when studying galaxy evolution, particularly with high-redshift observations from advanced observatories like the James Webb Space Telescope. Future investigations will explore the impact of AGN across a broader range of colour spaces and delve into the effects of intermediate-type AGN, paving the way for a more comprehensive understanding of the AGN population and its influence on galaxy formation and evolution studies. \\\\ \noindent \noindent \textbf{Keywords:} Active Galactic Nuclei, Supermassive Black Holes, Spectral Energy Distribution, SED Modelling, Photometry, Colour-Colour Diagrams, Colour Space, ZFOURGE
%\lipsum[1-6]